\documentclass[12pt,a4paper]{article}
\usepackage[utf8]{inputenc}
\usepackage[T1]{fontenc}
\usepackage[english,ngerman]{babel}
\usepackage{amsmath}
\usepackage{amsfonts}
\usepackage{amssymb}
\usepackage{booktabs}
\usepackage{geometry}
\usepackage{hyperref}
\usepackage{xcolor}

\geometry{margin=2.5cm}

\title{Dokumentation: Quantitative Research Data Infrastructure \\ \large Aufbau eines Financial Data Lake (PostgreSQL)}
\author{Research Infrastructure Project}
\date{\today}

\begin{document}

\maketitle

\section{Einleitung und Zielsetzung}
Das Ziel dieses Projekts ist der Aufbau einer robusten, skalierbaren und wissenschaftlich belastbaren Dateninfrastruktur zur Analyse von Finanzmarktdaten. Die Architektur ist so konzipiert, dass sie sowohl öffentliche Aktienmärkte (Public Equity) als auch Private Equity Deals, Derivate und unstrukturierte Textdaten (SEC Filings) unter einem Dach vereint.

Besonderer Wert wird auf folgende Punkte gelegt:
\begin{itemize}
    \item \textbf{Vermeidung von Survivorship Bias:} Speicherung historischer Index-Zusammensetzungen (Point-in-Time), um Backtesting-Fehler zu vermeiden.
    \item \textbf{Datenintegrität:} Einsatz von Relationalen Datenbanken (PostgreSQL) mit strikten Primärschlüsseln, um Dubletten und Inkonsistenzen zu verhindern.
    \item \textbf{Automatisierung:} Python-basierte Synchronisations-Skripte mit integrierter Abfrage-Logik zur Minimierung von API-Traffic (Refinitiv Eikon).
    \item \textbf{Sicherheit:} Strikte Trennung von Programmcode und Zugangsdaten über Umgebungsvariablen (\texttt{.env}).
\end{itemize}

\section{Datenbank-Architektur}
Um eine logische Trennung verschiedener Assetklassen bei gleichzeitiger Verknüpfbarkeit zu gewährleisten, wird eine Struktur basierend auf \textbf{Schemas} innerhalb einer zentralen Datenbank (\texttt{research\_db}) verwendet.

\subsection{Schema-Struktur}
\begin{table}[h]
\centering
\begin{tabular}{@{}llp{8cm}@{}}
\toprule
\textbf{Schema} & \textbf{Fokus} & \textbf{Beschreibung} \\ \midrule
\texttt{pub\_equity} & Public Equity & Börsennotierte Wertpapiere, Indizes, Preise und Optionen. \\
\texttt{priv\_equity} & Private Equity & Fondsdaten, Deal-Strukturen und Multiples. \\
\texttt{sec} & Regulatory & Rohdaten und Metadaten von SEC Filings (z.B. 10-K, 10-Q). \\
\texttt{master} & Global Meta & Sektor-Klassifizierungen und globale Entitäten. \\ \bottomrule
\end{tabular}
\end{table}

\section{Aktueller Implementierungsstand (Public Equity)}
Bisher wurde das Fundament für die Analyse von Aktienindizes am Beispiel des \textit{Dow Jones Industrial Average} (.DJI) gelegt.

\subsection{Tabelle: \texttt{index\_overview}}
Speichert die Stammdaten der beobachteten Indizes.
\begin{itemize}
    \item \texttt{index\_ric}: Eindeutiger Reuters Instrument Code (Primary Key).
    \item \texttt{index\_name}: Klarname des Index.
\end{itemize}

\subsection{Tabelle: \texttt{index\_constituents}}
Speichert die Zusammensetzung der Indizes zu diskreten Zeitpunkten (Snapshots).
\begin{itemize}
    \item \textbf{Logik:} Monatliche Snapshots (MS) ab Januar 2015.
    \item \textbf{Spalten:} \texttt{index\_ric}, \texttt{constituent\_ric}, \texttt{symbol}, \texttt{company\_name}, \texttt{snapshot\_date}.
    \item \textbf{Primary Key:} (\texttt{index\_ric}, \texttt{constituent\_ric}, \texttt{snapshot\_date}) stellt sicher, dass eine Aktie in mehreren Indizes gleichzeitig existieren kann, ohne Redundanz-Probleme zu erzeugen.
\end{itemize}

\section{Roadmap und geplante Erweiterungen}
Die nächsten Ausbaustufen konzentrieren sich auf die Anreicherung der bestehenden Listen mit Markt- und Fundamentaldaten.

\begin{enumerate}
    \item \textbf{Market Data Sync (\texttt{prices\_daily}):} Automatisierter Download von Total Return Preisen für alle in \texttt{index\_constituents} gelisteten Ticker.
    \item \textbf{Options Surface (\texttt{options\_chains}):} Speicherung von Optionsketten (Strikes, Expiries, IVs) zur Berechnung von Volatilitätsflächen.
    \item \textbf{SEC Crawler:} Verknüpfung der Ticker mit CIK-Nummern, um automatisiert Textdaten von der SEC-Website für Sentiment-Analysen zu ziehen.
    \item \textbf{Cross-Asset Analysis:} Integration der \texttt{priv\_equity} Daten, um Korrelationen zwischen Private Equity Bewertungen und Public Market Multiples zu untersuchen.
\end{enumerate}

\end{document}